\documentclass[11pt,a4paper]{article}

\usepackage{kotex}
\usepackage{fontspec}
\setmainfont{Times New Roman}
\setmainhangulfont{AppleMyungjo}
\setsanshangulfont{Apple SD Gothic Neo}

\usepackage[a4paper,top=18mm,bottom=18mm,left=20mm,right=20mm]{geometry}
\usepackage{fancyhdr}
\setlength{\headheight}{24pt}
\usepackage{enumitem}
\usepackage{mdframed}
\usepackage{array}
\usepackage{booktabs}

\setlength{\parindent}{1em}
\linespread{1.18}

\pagestyle{fancy}
\fancyhf{}
\renewcommand{\headrulewidth}{0.6pt}
\fancyhead[L]{\sffamily\bfseries 2026 고1 영어 변형 — 정답 및 해설 지문 27}
\fancyhead[R]{\sffamily [Junior Coding Camp 안내문]}
\renewcommand{\footrulewidth}{0pt}
\fancyfoot[C]{\framebox[16mm][c]{\thepage}}

\newmdenv[linewidth=0.6pt,roundcorner=3pt,innertopmargin=7pt,innerbottommargin=7pt,
          innerleftmargin=9pt,innerrightmargin=9pt,skipabove=5pt,skipbelow=5pt]{pbox}

\begin{document}

\begin{center}
{\fontsize{15}{17}\selectfont\sffamily\bfseries [지문 27] 정답 및 해설}
\end{center}
\vspace{6pt}

%% 정답 표
\begin{center}
\begin{tabular}{|c|c|c|c|}
\hline
\textbf{문항} & \textbf{유형} & \textbf{정답} & \textbf{배점} \\
\hline
1 & 내용 불일치 & ④ & 2점 \\
\hline
2-(1) & 서술형 단답 & 10세 이상 13세 이하 & 1점 \\
\hline
2-(2) & 서술형 단답 & on site & 1점 \\
\hline
3 & 오후 활동 & ② & 2점 \\
\hline
4 & 추론 & ③ & 2점 \\
\hline
\end{tabular}
\end{center}

\vspace{10pt}

%% 문항별 해설
{\sffamily\bfseries 1번 해설} (내용 불일치 — 2점)

\vspace{4pt}
\noindent \textbf{정답: ④}

\vspace{4pt}
\noindent \textbf{정답 근거:} 안내문에 ``Participation Fee: \$100 per person (\textbf{lunch provided})''라고 명시되어 있어,
참가비에 점심 식사가 포함된다. 따라서 ``참가자가 자신의 식사를 직접 가져와야 하며 음식이 제공되지 않는다''는 ④의 진술은 본문과 일치하지 않는다.

\vspace{6pt}
\noindent \textbf{오답 소거:}
\begin{itemize}[leftmargin=2em,itemsep=2pt,topsep=3pt,parsep=0pt]
  \item ①\ ``The event spans a single day at a technology facility.'' → 날짜가 7월 29일 하루이며 장소가 Arendale \textbf{Tech} Center이므로 일치. (TRUE)
  \item ②\ ``Participation is restricted to those with no prior coding experience.'' → ``for beginners''에 해당하므로 일치. (TRUE)
  \item ③\ ``Those who register early are given priority for spots.'' → ``On a first-come, first-served basis''와 일치. (TRUE)
  \item ⑤\ ``The morning session focuses on foundational coding skills.'' → 오전 9 a.m.--12 p.m.에 ``Coding Basics''가 배정되어 있으므로 일치. (TRUE)
\end{itemize}

\vspace{10pt}

{\sffamily\bfseries 2번 해설} (서술형 단답 — 각 1점)

\vspace{6pt}
\noindent \textbf{(1) 모범답안:} 10세 이상 13세 이하 (또는 ``10\textasciitilde{}13세'')

\vspace{3pt}
\noindent \textbf{채점 기준:}
\begin{itemize}[leftmargin=2em,itemsep=1pt,topsep=2pt,parsep=0pt]
  \item 10세와 13세가 모두 포함된 연령 범위를 정확히 서술한 경우 1점 부여.
  \item 한쪽 경계만 언급하거나 범위가 틀린 경우 0점.
\end{itemize}

\vspace{6pt}
\noindent \textbf{(2) 모범답안:} on site (또는 ``Register on site.'')

\vspace{3pt}
\noindent \textbf{채점 기준:}
\begin{itemize}[leftmargin=2em,itemsep=1pt,topsep=2pt,parsep=0pt]
  \item ``on site''가 포함된 경우 1점 부여.
  \item 철자 오류 없이 정확히 기재해야 함.
\end{itemize}

\vspace{10pt}
\noindent{\sffamily\bfseries 3번 해설} \quad \textbf{정답: ②}
\par 오후 1시부터 4시까지의 활동은 \textit{Game Design}이므로 ②가 정답이다.

\vspace{8pt}
\noindent{\sffamily\bfseries 4번 해설} \quad \textbf{정답: ③}
\par 안내문에 ``Limited to 50 students''와 ``On a first-come, first-served basis''라고 되어 있으므로, 먼저 등록한 학생이 참가 가능성이 높다고 추론할 수 있다.

\end{document}
